% resume-screen-two-page.tex — medium=screen (issue #184): a multi-page resume
% must carry no folio on any page and no running header on page two.
%
% The page break is forced rather than produced by overflowing content: this
% fixture asserts which furniture is emitted, not where the class chooses to
% break, and the sibling *-two-page fixtures already stress the latter under
% the default medium=print.
\documentclass[fontsize=11pt, margin=narrow, medium=screen]{careerdossier-resume}
\CDossierSetup{
  name     = {Ada Lovelace},
  headline = {Data Scientist},
  email    = {ada.lovelace@example.com},
  phone    = {+1 416 555 0142},
  location = {Toronto, Ontario, Canada},
}
\begin{document}
\MakeCDossierHeader

\CDossierSection{Summary}
Data scientist with a decade of experience across public-sector analytics,
reporting platforms, and reproducible research pipelines.

\CDossierSection{Experience}
\begin{CDossierEntry}[title = {Principal Data Scientist}, organization = {Digital Services Laboratory}, location = {Ottawa, ON}, dates = {2024--Present}]
  \begin{CDossierItemize}
    \item Led the delivery of reporting services supporting high-volume public programs and internal operational dashboards.
    \item Established automated validation for data quality, reproducibility, and archival output across every published dataset.
  \end{CDossierItemize}
\end{CDossierEntry}

\newpage

\CDossierSection{Education}
\begin{CDossierEntry}[title = {Ph.D. in Computer Science}, organization = {University of Toronto}, location = {Toronto, ON}, dates = {2021--2026}]
\end{CDossierEntry}
\begin{CDossierEntry}[title = {B.A.Sc. in Computer Engineering, Honours}, organization = {University of Waterloo}, location = {Waterloo, ON}, dates = {2014--2019}]
\end{CDossierEntry}

\end{document}
